\honest{If Many-Worlds Had Come First}{Eliezer Yudkowsky}{2008}

I do not claim I would have done better in that era.

Many-worlds was first proposed by Hugh Everett in 1957. The paper was ignored. John Wheeler sent Everett to see Niels Bohr, and Bohr did not take him seriously. ``Crushed, Everett left academic physics,'' invented ``the general use of Lagrange multipliers in optimization problems,'' and became a multimillionaire. \nb{Everett had taken a Pentagon job in 1956, to keep his student draft deferment, before the paper appeared and three years before he met Bohr. Wheeler later said he was ``disappointed, perhaps bitter, at the nonreaction to his theory.'' Lagrange multipliers were in use long before Everett; his 1963 paper gave a generalized method for allocating resources.} The field heard of his idea only in 1970, from Bryce DeWitt, who named it ``many-worlds.'' It may now be the majority view, ``(or not).'' \nb{In the polls available it has been a minority view.}

Suppose many-worlds had come first, in the 1920s, and no one had proposed a collapse until 1957. Would collapse now be gaining ground?

In my alternate world, Huve Erett brings Biels Nohr a one-sentence paper: at a measurement, all of the wavefunction except one point vanishes, the survivor chosen by the Born statistics. Nohr asks what a measurement is, and Huve changes it to ``When a quantum superposition gets too large.'' How large? Fifty microns, Huve guesses, and a student announces that superposition has just been verified at fifty microns. ``Whichever level, then.''

Nohr lists what the collapse would be: the only law in quantum mechanics that is non-linear, non-unitary and discontinuous, the only thing in physics with a preferred basis, the only thing that violates CPT symmetry, Liouville's Theorem and special relativity. \nb{The list follows the one Yudkowsky gave the day before, in ``Collapse Postulates.'' On relativity, a collapse model consistent with it was published in 2006, though only for particles that do not interact.} It explains the Born statistics no better than little angels would explain Maxwell's equations. It is ``post hoc.'' And it would be ``the only informally specified, qualitative fundamental law'' in physics. \nb{That fits Huve's postulate and the textbook one. It does not fit the GRW theory of 1986, which fixes where, how often and in which basis the collapse happens. The post does not mention it.} Worlds, Nohr says, are not fundamental; the fundamental law is the well-tested ``Wheeler-DeWitt equation.'' \nb{That equation is a proposed equation of quantum gravity. The tested one is the Schrödinger equation.} Huve asks for a test; Nohr says that without a specified collapse he cannot design one. \nb{Collapse theories that fix a rate can be tested. One, the simplest Diósi--Penrose model, was ruled out in 2021; none of the tests has found a collapse.}

Nohr admits that his side does not know why the Born statistics hold. \nb{The matching problem for many-worlds, how there can be probabilities when every outcome occurs, is ``arguably the most serious difficulty'' for it, according to the \textsc{Stanford Encyclopedia}. Huve does not raise it.} Huve says the collapse cannot be used to signal faster than light, and then that the relevant properties ``don't exist'' until the collapse. \nb{Two days earlier, in ``Quantum Non-Realism,'' Yudkowsky had attacked that move.} Nohr tells him to ``Just say `oops'.''

At the end Huve imagines a world where collapse came first and, for thirty years, no one suggested many worlds. Nohr laughs. Physicists who test equations in every case they can, and never ask whether the equations are universal? Von Neumann, ``who may have been the highest-g human in history,'' would not think of it? Schrödinger would carry the superposition up to a cat, then stop and say he could not imagine what happens next? \nb{In our world von Neumann applied the Schrödinger equation to the measuring instrument in 1932 and kept the collapse at the observer. Schrödinger, in 1935, did carry the superposition up to the cat, called it one of the ``quite ridiculous cases,'' and said this prevents us from naively accepting a ``blurred model'' of reality. They asked the question and judged the answer absurd.}

Nohr's last line: this would be a world where everyone is ``too dumb to sign up for cryonics,'' with George W. Bush as president.
